\chapter{Apprentissage automatique}\label{ch:iv1:machine-learning}

\begin{ivtrains}
L'apprentissage automatique tel que les entretiens le testent~: biais et variance, arbres et boosting, fuites d'information et validation sur séries temporelles, régularisation, réseaux de neurones à la main, métriques, et ce qu'un signal faible fait à tout cela. Presque chaque réponse ici est un nombre que l'on peut dériver.
\end{ivtrains}

\begin{question}[$\star$ \iqroles{QR, DS, Analyste} \iqfirm{toutes} \iqlevel{junior}]\label{iq:iv1:machine-learning:1}
Un modèle signale 100 jours sur 1000 comme des jours de forte hausse. Parmi les jours signalés, 40 en étaient réellement, et 10 autres jours de ce type n'ont pas été signalés. Donnez l'exactitude, la précision, le rappel et le F1~: quelle exactitude atteint un modèle qui ne signale jamais rien~?
\end{question}

\begin{question}[$\star$ \iqroles{QR, DS, QT} \iqfirm{toutes} \iqlevel{junior}]\label{iq:iv1:machine-learning:2}
Vous essayez 20 modèles sans aucun pouvoir prédictif, chacun sur le même ensemble de test de 100 jours indépendants à la hausse ou à la baisse, les modèles étant indépendants les uns des autres, et vous présentez le meilleur. Quelle exactitude annoncez-vous à peu près~? Qu'auriez-vous dû faire autrement~?
\end{question}

\begin{question}[$\star$ \iqroles{QR, DS, QD} \iqfirm{toutes} \iqlevel{junior}]\label{iq:iv1:machine-learning:3}
Vous enregistrez une variable en cents plutôt qu'en dollars, puis vous réajustez les modèles. Lesquels de ces modèles changent leurs prédictions~: moindres carrés ordinaires, régression ridge, $k$ plus proches voisins, un arbre de décision, une forêt aléatoire, une régression sur les deux premières composantes principales~?
\end{question}

\begin{question}[$\star$ \iqroles{QR, DS, Analyste} \iqfirm{toutes} \iqlevel{junior}]\label{iq:iv1:machine-learning:4}
Une étiquette binaire vaut 1 sur 10~\% des lignes. Quelle log loss obtient la prédiction constante $0.1$, et la constante $0.5$~? Quelle perte sur une seule ligne résulte de la prédiction $0.999$ sur une ligne dont l'étiquette vaut 0, et pourquoi cela doit-il vous inquiéter~?
\end{question}

\begin{question}[$\star$ \iqroles{QR, DS, Analyste} \iqfirm{toutes} \iqlevel{junior}]\label{iq:iv1:machine-learning:5}
Une régression logistique affiche un coefficient de $0.7$ pour une variable standardisée. Que deviennent les cotes lorsque la variable augmente de 1~? En partant des probabilités $0.2$ puis $0.5$, quelle est la nouvelle probabilité dans chaque cas~?
\end{question}

\begin{question}[$\star$ \iqroles{QR, DS, QD} \iqfirm{toutes} \iqlevel{junior}]\label{iq:iv1:machine-learning:6}
Vous lancez une validation croisée à 5 plis sur 1000 lignes, sur une grille de 10 valeurs pour chacun des 3 hyperparamètres. Quelle est la taille de chaque ensemble d'entraînement~? Combien de fois une ligne donnée est-elle validée, et à combien de modèles aide-t-elle à entraîner~? Combien d'ajustements la recherche requiert-elle au total~?
\end{question}

\begin{question}[$\star$ \iqroles{QR, DS, Risque} \iqfirm{toutes} \iqlevel{junior}]\label{iq:iv1:machine-learning:7}
Les mauvaises transactions représentent 1~\% d'un flux. Un détecteur en repère 90~\% et signale à tort 5~\% des bonnes. Quelle part des signalements est réelle~? Quelle serait cette part si le taux de faux positifs tombait à 1~\%~? L'aire sous la courbe ROC change-t-elle si les mauvaises transactions passent à 10~\% du flux~?
\end{question}

\begin{question}[$\star$ \iqroles{QR, DS, QD} \iqfirm{toutes} \iqlevel{junior}]\label{iq:iv1:machine-learning:8}
Un réseau entièrement connecté a 20 entrées, des couches cachées de 64 et 32 unités et une sortie, avec un biais pour chaque unité qui calcule. Combien de paramètres possède-t-il~? Vous l'entraînez sur 1250 lignes quotidiennes. À quoi cela doit-il vous faire penser~?
\end{question}

\begin{question}[$\star$ \iqroles{QR, DS, QT} \iqfirm{toutes} \iqlevel{junior}]\label{iq:iv1:machine-learning:9}
Deux modèles, visant une exactitude cible de 85~\%. Le modèle A obtient 99~\% à l'entraînement et 70~\% en validation. Le modèle B obtient 72~\% et 70~\%. Quel est le problème probable de chacun~? Pour chacun, des données supplémentaires aident-elles, et que modifieriez-vous~?
\end{question}

\begin{question}[$\star\star$ \iqroles{QR, DS, QD} \iqfirm{toutes} \iqlevel{junior}]\label{iq:iv1:machine-learning:10}
La vraie fonction est $f(x) = x$, et les observations ont un bruit de variance 1. Vous estimez $f$ en un point à l'aide de la moyenne de ses $k$ plus proches voisins, situés aux distances $0.2, 0.4, \dots, 0.2k$, tous d'un même côté. Écrivez le biais et la variance. Quelle valeur de $k$ minimise l'erreur quadratique de l'estimation~? Que devient le meilleur $k$ si les données sont deux fois plus denses~?
\end{question}

\begin{question}[$\star\star$ \iqroles{QR, DS} \iqfirm{toutes} \iqlevel{junior}]\label{iq:iv1:machine-learning:11}
Un modèle linéaire correct à 20 paramètres est ajusté sur 100 lignes, avec une variance du bruit égale à 1 et des régresseurs fixes. Quelle est l'erreur quadratique moyenne espérée sur les lignes d'entraînement~? Quelle est-elle sur un nouveau bruit aux mêmes régresseurs~? Interprétez l'écart, et dites comment il dépend du nombre de paramètres.
\end{question}

\begin{question}[$\star\star$ \iqroles{QR, DS, Analyste} \iqfirm{toutes} \iqlevel{junior}]\label{iq:iv1:machine-learning:12}
Un nœud compte 40 positifs et 60 négatifs. Une division envoie 30 positifs et 10 négatifs à gauche, et le reste à droite. Calculer la diminution de l'impureté de Gini et celle de l'entropie, en bits.
\end{question}

\begin{question}[$\star\star$ \iqroles{QR, DS} \iqfirm{toutes} \iqlevel{junior}]\label{iq:iv1:machine-learning:13}
Un nœud compte 400 positifs et 400 négatifs. La division A donne les enfants (300, 100) et (100, 300). La division B donne (200, 400) et (200, 0). Comparez-les par l'erreur de classification et par l'impureté de Gini. Laquelle préférez-vous, et quelle propriété générale du critère explique le résultat~?
\end{question}

\begin{question}[$\star\star$ \iqroles{QR, DS, QT} \iqfirm{toutes} \iqlevel{junior, confirmé}]\label{iq:iv1:machine-learning:14}
Vous moyennez $B$ modèles, chacun de variance $\sigma^2$, et de corrélation deux à deux $\rho$. Dérivez la variance de la moyenne. Pour $\sigma^2 = 1$, $\rho = 0.3$ et $B = 100$, quelle vaut-elle, et que devient-elle quand $B$ croît~? En utilisant à chaque division des sous-ensembles aléatoires de variables, on abaisse $\rho$ à $0.1$ mais on fait monter $\sigma^2$ à $1.2$. Est-ce un bon compromis~?
\end{question}

\begin{question}[$\star\star$ \iqroles{QR, DS} \iqfirm{toutes} \iqlevel{junior}]\label{iq:iv1:machine-learning:15}
Dans le gradient boosting à perte quadratique, montrez que le résidu est l'opposé du gradient. Supposez que chaque nouvel arbre ajuste exactement le résidu courant d'une ligne donnée, et qu'il soit ajouté avec un rétrécissement $\nu = 0.1$. Quelle part de la cible est ajustée après $m$ arbres~? Quand dépasse-t-on 90~\% puis 99~\%~? Pourquoi rétrécir du tout~?
\end{question}

\begin{question}[$\star\star$ \iqroles{QR, DS, QT} \iqfirm{hedge fund} \iqlevel{junior, confirmé}]\label{iq:iv1:machine-learning:16}
Les données présentent un seul motif réel dans deux variables, et 20~\% des étiquettes sont inversées au hasard. Vous augmentez le nombre d'arbres à 1000, d'abord dans une forêt aléatoire d'arbres profonds, puis dans un boosting d'arbres petits avec un grand taux d'apprentissage. Comment évoluent les erreurs d'entraînement et de test dans chaque cas~? Pourquoi cette différence~?
\end{question}

\begin{question}[$\star\star$ \iqroles{QR, DS} \iqfirm{toutes} \iqlevel{junior}]\label{iq:iv1:machine-learning:17}
Dans AdaBoost, un apprenant faible a une erreur pondérée de $0.3$. Quel est son poids de vote $\alpha = \frac12\ln((1 - \varepsilon)/\varepsilon)$~? Par quels facteurs sont multipliés les poids des lignes mal classées et ceux des lignes bien classées, et quel poids total détiennent les lignes mal classées après normalisation~? Que signifie ce total pour l'apprenant suivant~?
\end{question}

\begin{question}[$\star\star$ \iqroles{QR, DS} \iqfirm{hedge fund} \iqlevel{junior, confirmé}]\label{iq:iv1:machine-learning:18}
Montrez que la régression ridge est la moyenne a posteriori sous un a priori gaussien, avec $\beta \sim N(0, \tau^2 I)$ et une variance du bruit $\sigma^2$, et déterminez $\lambda$. Définissez ensuite son nombre effectif de paramètres, et calculez-le pour des carrés de valeurs singulières de 100, 10 et 1, avec $\lambda = 10$. Vérifiez les deux limites de $\lambda$.
\end{question}

\begin{question}[$\star\star$ \iqroles{QR, DS} \iqfirm{toutes} \iqlevel{junior, confirmé}]\label{iq:iv1:machine-learning:19}
Une régression logistique sans terme constant, à une seule variable, comporte deux lignes~: $x = -1$ avec l'étiquette 0, et $x = 1$ avec l'étiquette 1. Que fait l'estimateur du maximum de vraisemblance~? Ajoutez une pénalité $\frac12\lambda\beta^2$ avec $\lambda = 0.1$. Quelle est désormais $\beta$, et quelle probabilité le modèle donne-t-il en $x = 1$~?
\end{question}

\begin{question}[$\star\star$ \iqroles{QR, DS} \iqfirm{toutes} \iqlevel{junior}]\label{iq:iv1:machine-learning:20}
Un classifieur bayésien naïf part d'un a priori de $\frac12$. La variable A a un rapport de vraisemblance de 3 pour la classe 1. Elle a été incluse par erreur trois fois. La variable B, incluse une seule fois, a un rapport de vraisemblance de $0.1$. Donnez la probabilité a posteriori vraie et celle du classifieur si A et B sont réellement indépendantes conditionnellement à la classe. La décision change-t-elle au seuil de $\frac12$~?
\end{question}

\begin{question}[$\star\star$ \iqroles{QR, DS, QD} \iqfirm{toutes} \iqlevel{junior}]\label{iq:iv1:machine-learning:21}
Un réseau a une seule entrée $x = 1$, une unité cachée sigmoïde $h = \sigma(w_1x)$, et une sortie $\hat{y} = w_2h$, avec $w_1 = 0$ et $w_2 = 2$. La cible est $y = 0$ et la perte $\frac12(\hat{y} - y)^2$. Calculez les deux gradients par la règle de la chaîne. Faites un pas de taille $0.1$ sur les deux poids, puis donnez la nouvelle perte.
\end{question}

\begin{question}[$\star\star$ \iqroles{QR, DS, QD} \iqfirm{toutes} \iqlevel{junior, confirmé}]\label{iq:iv1:machine-learning:22}
Dix couches sigmoïdes sont enchaînées, avec des poids égaux à 1. Quelle borne supérieure s'applique au gradient qui atteint la première couche~? Qu'est-ce que ReLU change, et quelle nouvelle défaillance cela introduit-il~? Dans une couche ReLU de largeur $n$, quelle variance des poids initiaux maintient la variance du signal constante à travers la profondeur~?
\end{question}

\begin{question}[$\star\star$ \iqroles{QR, DS, QD} \iqfirm{toutes} \iqlevel{junior}]\label{iq:iv1:machine-learning:23}
Démontrez qu'aucun classifieur linéaire ne calcule le ou exclusif de deux bits. Écrivez ensuite un réseau à une couche cachée de deux unités ReLU qui le calcule exactement, puis vérifiez les quatre entrées.
\end{question}

\begin{question}[$\star\star$ \iqroles{QR, DS, QD} \iqfirm{toutes} \iqlevel{junior}]\label{iq:iv1:machine-learning:24}
Une sortie sigmoïde a un logit $z = -7$, et l'étiquette vaut 1. Calculez le gradient par rapport à $z$ sous l'entropie croisée, puis sous la perte quadratique $\frac12(p - y)^2$. De quel facteur diffèrent-ils~? Chaque perte est-elle convexe en $z$~? Qu'est-ce que cela implique pour l'entraînement~?
\end{question}

\begin{question}[$\star\star$ \iqroles{QR, DS, Analyste} \iqfirm{toutes} \iqlevel{junior}]\label{iq:iv1:machine-learning:25}
Montrez que l'aire sous la courbe ROC égale la probabilité qu'un positif tiré au hasard obtienne un score supérieur à celui d'un négatif tiré au hasard, les ex æquo comptant pour moitié. Les positifs obtiennent des scores $0.9$, $0.6$, $0.4$, et les négatifs des scores $0.7$, $0.3$, $0.2$, $0.1$. Calculez l'aire par comptage des paires, puis par somme des rangs.
\end{question}

\begin{question}[$\star\star$ \iqroles{QR, DS, QT} \iqfirm{toutes} \iqlevel{junior}]\label{iq:iv1:machine-learning:26}
Une fausse alerte coûte 1 et une non-détection coûte 9. Au-dessus de quelle probabilité calibrée faut-il agir~? Maintenant, une transaction gagne $0.5~\%$ si elle a raison et perd 1~\% si elle a tort. Au-dessus de quelle probabilité vaut-il la peine de la prendre~? Que peut-il mal se passer si les scores ne sont pas calibrés~?
\end{question}

\begin{question}[$\star\star$ \iqroles{QR, QT, DS} \iqfirm{hedge fund, teneur de marché} \iqlevel{junior}]\label{iq:iv1:machine-learning:27}
Un modèle directionnel a raison dans 52~\% des jours, et le mouvement de chaque jour a la même taille, à savoir 1 unité. Quels ratios de Sharpe quotidien et annuel obtient une position quotidienne, en prenant 252 jours par an et sans coûts~? Et si la transaction de chaque jour coûte $0.02$ unité~? Quel coût par transaction ramène le ratio de Sharpe à zéro~?
\end{question}

\begin{question}[$\star\star$ \iqroles{QR, DS, QT} \iqfirm{hedge fund} \iqlevel{junior, confirmé}]\label{iq:iv1:machine-learning:28}
Votre étiquette de chaque jour est le rendement total des 20 jours suivants, et vous disposez de 1000 jours. Quelle est la corrélation entre les étiquettes séparées de $k$ jours, sous des rendements quotidiens indépendants~? Pour un bloc de test de 200 jours au milieu, combien de lignes d'entraînement faut-il retirer, pour qu'aucune étiquette d'entraînement ne partage un jour avec une étiquette de test~? Supposez que la seule variable soit la date et que le modèle soit le 1-plus proche voisin. Que rapporte une découpe aléatoire en 5 plis sur du pur bruit~?
\end{question}

\begin{question}[$\star\star$ \iqroles{QR, DS, QT} \iqfirm{hedge fund} \iqlevel{junior, confirmé}]\label{iq:iv1:machine-learning:29}
Un modèle prédit le rendement de chaque action entre l'ouverture et la clôture du jour $t$. Lesquels de ces éléments introduisent une fuite d'information, et comment~? (a) Le rendement clôture-à-clôture du jour $t - 1$. (b) Le secteur de l'action, selon la classification d'aujourd'hui, sur un historique de dix ans. (c) Des variables standardisées avec la moyenne et la variance de l'échantillon entier. (d) Les résultats trimestriels placés le dernier jour du trimestre de l'exercice. (e) Un indicateur indiquant si l'action appartient à l'indice, d'après les membres d'aujourd'hui. (f) Le volume du jour $t$.
\end{question}

\begin{question}[$\star\star$ \iqroles{QR, DS} \iqfirm{toutes} \iqlevel{junior}]\label{iq:iv1:machine-learning:30}
L'algorithme des $k$-moyennes est lancé sur les quatre points $(\pm 2, \pm 0.5)$, avec $k = 2$. Initialisez les centres en $(0, 0.5)$ et $(0, -0.5)$. Que se passe-t-il, et quelle est la somme des carrés~? Quel est l'optimum~? Si l'initialisation consiste en deux points de données différents tirés au hasard, quelle est la probabilité du mauvais résultat~? Avec cinq réinitialisations indépendantes~?
\end{question}

\begin{question}[$\star\star\star$ \iqroles{QR, DS} \iqfirm{hedge fund} \iqlevel{junior, confirmé}]\label{iq:iv1:machine-learning:31}
Vous disposez de 50 lignes, d'étiquettes aléatoires équilibrées et de 5000 variables de pur bruit. Vous conservez les 100 variables les plus corrélées aux étiquettes sur l'ensemble des données, puis vous lancez une validation croisée à 5 plis d'un classifieur au centroïde le plus proche sur ces 100 variables. Quelle erreur cela rapporte-t-il, et quelle est la vérité~? Donnez la bonne procédure, et montrez la différence par une simulation.
\end{question}

\begin{question}[$\star\star\star$ \iqroles{QR, QT, DS} \iqfirm{hedge fund} \iqlevel{confirmé}]\label{iq:iv1:machine-learning:32}
Les ratios de Sharpe annuels vrais d'un ensemble de signaux candidats suivent des lois normales centrées en 0, d'écart type $0.5$, et un backtest mesure chacun avec un bruit d'écart type 1. Vous en testez 100 de façon indépendante et sélectionnez le meilleur. Quel ratio de Sharpe vous attendez-vous à mesurer sur celui-ci, et quelle est son espérance de ratio de Sharpe vrai~? Comment les chiffres changent-ils avec quatre fois plus d'historique~?
\end{question}

\begin{question}[$\star\star\star$ \iqroles{QR, DS, QT} \iqfirm{hedge fund} \iqlevel{junior, confirmé}]\label{iq:iv1:machine-learning:33}
Une moyenne qui évolue lentement suit une marche aléatoire d'innovation de variance $q = 0.0001$, et vous l'observez avec un bruit de variance 1. Vous prévoyez sa valeur suivante par la moyenne des $w$ dernières observations. Dérivez l'erreur quadratique moyenne de la prévision en fonction de $w$, ainsi que le meilleur $w$. Dans quelle mesure la réponse est-elle sensible à $w$~? Qu'est-ce que cela dit du choix d'une fenêtre d'entraînement~?
\end{question}

\begin{question}[$\star\star\star$ \iqroles{QR, DS, QT} \iqfirm{hedge fund} \iqlevel{junior, confirmé}]\label{iq:iv1:machine-learning:34}
Les scores suivent des lois normales de variance 1, de moyenne 0 pour les négatifs et de moyenne $d$ pour les positifs. Dérivez l'aire sous la courbe ROC. Quelles sont les aires pour $d = 0.1$, $0.5$ et 1~? Un modèle de rendement rapporte une aire de $0.55$ sur des classes équilibrées. Quelle est la valeur de $d$, et quelle corrélation entre le score et l'étiquette cela implique-t-il~?
\end{question}

\begin{question}[$\star\star\star$ \iqroles{QR, DS, Risque} \iqfirm{toutes} \iqlevel{junior, confirmé}]\label{iq:iv1:machine-learning:35}
Vous conservez tous les positifs mais seulement 10~\% des négatifs, tirés au hasard, puis vous entraînez un classifieur probabiliste. Il renvoie $0.5$ et $0.9$ sur deux lignes. Quelles sont les probabilités correctes pour la population d'origine~? Qu'est-ce qui change dans une régression logistique, et qu'est-ce qui ne change pas~? Dérivez la formule.
\end{question}

\begin{question}[$\star\star\star$ \iqroles{QR, DS} \iqfirm{toutes} \iqlevel{junior, confirmé}]\label{iq:iv1:machine-learning:36}
Le gradient boosting avec perte logarithmique démarre sur un jeu de données de 100 observations dont 60 positives. Déterminez le logit initial, ainsi que le gradient et la hessienne de chaque ligne. Un arbre isole une feuille de 30 positifs et 10 négatifs. Déterminez la valeur de feuille de Newton, et comparez-la à l'optimum exact de cette feuille. Avec un rétrécissement de $0.1$, quelle est la nouvelle probabilité de la feuille~?
\end{question}

\begin{question}[$\star\star\star$ \iqroles{QR, DS} \iqfirm{toutes} \iqlevel{junior, confirmé}]\label{iq:iv1:machine-learning:37}
Une bibliothèque de boosting du second ordre approxime la perte dans une feuille par $Gw + \frac12(H + \lambda)w^2$, où $G$ et $H$ sont les sommes des gradients et des hessiennes, et $w$ est la valeur de la feuille. Dérivez le meilleur $w$ et le gain d'une division. Comparez, avec $\lambda = 1$, la division $(G_L, H_L) = (-10, 20)$ et $(G_R, H_R) = (10, 20)$ à la division $(-10, 1)$ et $(10, 39)$. Laquelle des deux divisions la bibliothèque privilégie-t-elle, et quel réglage protège contre ce choix~?
\end{question}

\begin{question}[$\star\star\star$ \iqroles{QR, DS, QD} \iqfirm{toutes} \iqlevel{junior, confirmé}]\label{iq:iv1:machine-learning:38}
Prouvez que, pour les moindres carrés, le résidu leave-one-out vaut $e_i/(1 - h_{ii})$, où $e_i$ est le résidu ordinaire et $h_{ii}$ le levier. Ajustez ensuite une droite avec ordonnée à l'origine aux données $x = 0, 1, 2, 3, 10$ et $y = 0, 1, 2, 3, 0$, et calculez le résidu leave-one-out du dernier point à la fois par la formule et à partir de sa signification.
\end{question}

\begin{question}[$\star\star\star$ \iqroles{QR, DS} \iqfirm{toutes} \iqlevel{junior, confirmé}]\label{iq:iv1:machine-learning:39}
En régression linéaire, chaque entrée est multipliée au hasard par $\delta_j/q$, où $\delta_j$ vaut 1 avec probabilité $q$ et 0 sinon, ce qui constitue une forme de dropout. Montrez que l'espérance de la perte quadratique est une pénalité ridge sur les coefficients, et déterminez ses poids. Pour des variables non corrélées avec $\sum_i x_{ij}^2 = n$, comment les coefficients se comparent-ils aux moindres carrés ordinaires pour $q = 0.5$ et $q = 0.8$~?
\end{question}

\begin{question}[$\star\star\star$ \iqroles{QR, DS, QD} \iqfirm{toutes} \iqlevel{junior, confirmé}]\label{iq:iv1:machine-learning:40}
Pour une entrée sur $[0, 1]$, un réseau ReLU à une seule couche cachée et $n$ unités donne au plus combien de morceaux linéaires~? Écrivez l'application tente, $T(x) = 2x$ sur la première moitié et $2 - 2x$ sur la seconde, avec deux unités ReLU. Combien de morceaux possède l'application composée avec elle-même $d$ fois~? Comparez le nombre d'unités nécessaires au réseau profond et au réseau peu profond pour $d = 10$.
\end{question}

\begin{question}[$\star\star\star$ \iqroles{QR, DS} \iqfirm{hedge fund} \iqlevel{junior, confirmé}]\label{iq:iv1:machine-learning:41}
Pour le lasso, avec l'objectif $\frac{1}{2n}\|y - X\beta\|^2 + \lambda\|\beta\|_1$ et des colonnes standardisées, montrez que tous les coefficients sont nuls si et seulement si $\lambda \ge \max_j|x_j^\top y|/n$. Prenez $n = 1000$, $p = 1000$, des colonnes indépendantes et un $y$ de pur bruit de variance 1. À quoi correspond ce seuil~? Que dit-il du choix de $\lambda$ lorsque le signal est faible~?
\end{question}

\begin{question}[$\star\star\star$ \iqroles{QR, DS, QT} \iqfirm{hedge fund} \iqlevel{junior, confirmé}]\label{iq:iv1:machine-learning:42}
Une forêt aléatoire utilise une variable réelle, une copie exacte de celle-ci, et une variable inutile dont la valeur diffère à chaque ligne. Que disent l'importance par impureté et l'importance par permutation calculée hors échantillon au sujet de la variable réelle et de sa copie~? Et qu'en est-il de la variable inutile~? Comment évalueriez-vous un groupe de variables corrélées~?
\end{question}

\begin{question}[$\star\star\star\star$ \iqroles{QR, DS} \iqfirm{hedge fund} \iqlevel{confirmé}]\label{iq:iv1:machine-learning:43}
Régressez du pur bruit de variance 1, $y$ indépendant des variables, sur $p$ variables gaussiennes centrées réduites indépendantes, avec $n = 100$ lignes et les moindres carrés, en prenant la solution de norme minimale lorsque $p > n$. Dérivez l'erreur de test espérée sur une nouvelle ligne en fonction de $p$, pour $p$ inférieur et supérieur à $n$. Évaluez-la pour $p = 50, 90, 110, 200, 1000$. Pourquoi atteint-elle un maximum, puis pourquoi redescend-elle~?
\end{question}

\begin{question}[$\star\star\star\star$ \iqroles{QR, DS} \iqfirm{hedge fund} \iqlevel{confirmé}]\label{iq:iv1:machine-learning:44}
Deux classes gaussiennes de covariance identité en $p$ dimensions ont des moyennes qui diffèrent d'un vecteur de longueur 2 situé uniquement sur la première coordonnée, de sorte que $p - 1$ variables sont inutiles. Les deux classes sont équiprobables. Vous estimez les moyennes à partir de 50 lignes par classe, et vous classez à l'aide de l'hyperplan à mi-chemin entre les moyennes estimées. Dérivez l'erreur de test approchée en fonction de $p$, et évaluez-la pour $p = 1$, 100 et 1000. Quelle est l'erreur avec les moyennes connues~?
\end{question}

\begin{question}[$\star\star\star\star$ \iqroles{QR, DS} \iqfirm{hedge fund} \iqlevel{confirmé}]\label{iq:iv1:machine-learning:45}
Prouvez que l'erreur d'entraînement d'AdaBoost est au plus $\prod_t 2\sqrt{\varepsilon_t(1 - \varepsilon_t)}$, où $\varepsilon_t$ est l'erreur pondérée au tour $t$, et que celle-ci est au plus $\exp(-2\sum_t\gamma_t^2)$, avec $\gamma_t = \frac12 - \varepsilon_t$. Si chaque tour a $\varepsilon = 0.4$, combien de tours chacune des deux bornes exige-t-elle pour atteindre une erreur d'entraînement inférieure à 1~\%~? Pourquoi la borne ne requiert-elle aucune hypothèse sur les données~?
\end{question}
